%%
%% fall-lecture20.tex
%% 
%% Made by Alex Nelson <pqnelson@gmail.com>
%% Login   <alex@lisp>
%% 
%% Started on  2026-11-12T13:13:52-0800
%% Last update 2026-11-12T13:13:52-0800
%% 

\lecture{}

\begin{definition}
Let $\omega$ be a complex rank-$n$ bundle (and $\omega_{\RR}$ its
associated real rank-$2n$ bundle, with Euler class $e_{\RR}\in H^{2n}(B;\ZZ)$).
The \define{Top Chern Class} is $c_{n}(\omega):=e(\omega_{\RR})\in H^{2n}(B;\ZZ)$.
\end{definition}

\begin{node}
We need to introduce the Gysin sequence. First, we note the long exact
sequence of homology groups
\begin{equation}
\dots\gets H_{j-1}(E_{0})\gets H_{j}(E,E_{0})\gets H_{j}(E)\gets H_{j}(E_{0})\gets\dots
\end{equation}
so the long exact sequence of cohomology groups
\begin{equation}
  \dots\to H^{j}(E,E_{0})\xrightarrow{j}H^{j}(E)\xrightarrow{r}H^{j}(E_{0}\xrightarrow{\delta}H^{j+1}(E,E_{0})\to\dots
\end{equation}
where $r$ is the restriction morphism, $\delta$ is the connecting
morphism. We use the Thomas isomorphism $\varphi\colon H^{j-2n}(B)\to H^{j}(E,E_{0})$
and $\pi^{*}\colon H^{j}(B)\to H^{j}(E)$ to obtain the commutative
diagram
\begin{equation}
\vcenter{\xymatrix{H^{j}(E,E_{0})\ar[r]^{j} & H^{j}(E)\ar[r]^{r}&H^{j}(E_{0})\\
H^{j-2n}(B)\ar[u]^{\varphi}\ar@{..>}[r]&H^{j}(B)\ar[u]^{\pi^{*}}\ar@{..>}[ur]_{\pi_{0}^{*}}&\\
\beta\ar@{|->}[r] & (\pi^{*})^{-1}j(\pi^{*}\beta\smile U) & }}
\end{equation}
We see that
\begin{subequations}
  \begin{align}
(\pi^{*})^{-1}j(\pi^{*}\beta\smile U) &= (\pi^{*})^{-1}(\pi^{*}\beta\smile(\pi^{*})^{-1}j\pi^{*}U)\\
&= \beta\smile e(\omega_{\RR})
  \end{align}
\end{subequations}
This gives us a long exact sequence with terms like
\begin{equation}
\dots\to H^{j-2n}(B;\ZZ)\xrightarrow{\smile e(\omega_{\RR})}H^{j}(B)\xrightarrow{\pi_{0}^{*}}H^{j}(E_{0})\xrightarrow{\varphi}H^{j+1+2n}(B)\to\dots
\end{equation}
which is precisely the \define{Gysin Sequences}.
\end{node}

\begin{node}
We need a bundle $\omega_{0}$ over $E_{0}$ (to define the Chern
characters), where $\omega_{0}$ is a rank-$(n-1)$ complex bundle over $E_{0}$.
The basic idea is, just as $\pi\colon E\to B$ sends $(b,v)\mapsto b$,
we can consider the subspace of each fiber $v^{\perp}$ which is
($n-1$)-dimensional. This gives us a way to construct $\omega_{0}$.
\end{node}

\begin{definition}
We define the \define{Chern Character} on $\omega$ as
$c_{n}(\omega)=e(\omega_{\RR})$ [as before], and then inductively
$c_{n-1}(\omega)=(\pi_{0}^{*})^{-1}c_{n-1}(\omega_{0})$, $c_{n-2}(\omega)=(\pi_{0}^{*})^{-1}c_{n-2}(\omega_{0})=(\pi_{0}^{*})^{-1}(\pi_{00}^{*})^{-1}c_{n-2}((\omega_{0})_{0})$,
\dots, all the way down to $c_{0}(\omega)$.
\end{definition}

\begin{node}
There is a more intuitive way to think of this construction, but we
will discuss it next class.
\end{node}


\begin{definition}
The \define{Total Chern Class} is $c(\omega)=1+c_{1}(\omega)+\dots+c_{n}(\omega)$.
\end{definition}
\begin{proposition}
\begin{enumerate}
\item Naturality: for any $f\colon A\to B$ continuous, $c(f^{*}\omega)=f^{*}c(\omega)$
\item Product: $c(\omega_{1}\oplus\omega_{2})=c(\omega_{1})c(\omega_{2})$
\item $c_{1}(\gamma^{1}_{1})\in H^{2}(\CP^{1};\ZZ)$ is nonzero where
  $\gamma^{1}_{1}$ is the tautological line bundle over $\CP^{1}$.
\end{enumerate}
\end{proposition}

(We just need to show $c_{k}(\omega)=0$ for all $k>n$ when $\omega$ is
a rank-$n$ complex bundle, and then we'll have matched all the axioms
for the Stiefel--Whitney classes.)